% =========================
% puzzles/pXX.tex
% Final - The Clockwork Alibi Ledger
% =========================

\PuzzleChapterIntro{The Clockwork Alibi Ledger}{Temporal Logic \textbullet\ Scheduling \textbullet\ Alibi Verification}{This is a rigorous timeline puzzle. You must cross-reference travel times, mandatory wait periods, and mechanical timestamps to construct a minute-by-minute schedule for four custodians. The solution requires identifying the only suspect whose path is mathematically compatible with the crime window.}

\Lore{

The Ledger captures time. It is a machine of brass and pressure, indifferent to testimony.
At every threshold, a signet-plate demands a mark. Strike it, and the minute is burned into the record. Pass through, and the gate may hold you, the locking pins refusing to retract until the corridor decides you have waited long enough.

The Vault hums with internal rhythms: the Scriptorium press, the Furnace bell, the Gatehouse latch. These side-strips chatter only when a hand is present to work them. And deep in the center, the Reliquary door is timed by a spring that knows no mercy—it opens once, waits, and slams shut.
Four custodians swear innocence. The clockwork swears otherwise.

}

\Problem

The Vault has four relevant places: the \emph{Scriptorium} \(S\), the \emph{Furnace} \(F\), the \emph{Gatehouse} \(G\), and the \emph{Reliquary} \(R\).

\medskip

\noindent\textbf{Timing convention.}
All timestamps are exact \emph{to the minute}. Travel times and gate delays are in whole minutes.
You may wait anywhere, but you may not arrive earlier than the stated travel minima.

\medskip

\noindent\textbf{Travel times (door-to-door minima).}
\begin{itemize}
\item Each Way: | | \(G \leftrightarrow S\): \(5\) minutes | | \(G \leftrightarrow F\): \(5\) minutes | | \(S \leftrightarrow F\): \(4\) minutes | |
\item \emph{Whisper Stair} (down/up): \(S\to R\) takes \(4\) minutes, \(R\to S\) takes \(6\).
\item \emph{Ash Duct} (down/up): \(F\to R\) takes \(3\) minutes, \(R\to F\) takes \(5\).
\item Main hall: \(G \leftrightarrow R\) takes \(7\) minutes each way.
\end{itemize}

\medskip

\noindent\textbf{How the Ledger works.}
To \emph{enter} \(S\), \(F\), or \(G\), a custodian must press their own signet on the outside plate; the Ledger prints the minute of entry.
Exiting does \emph{not} print anything. The Gatehouse has no such delay.

Two wing-gates have pressure seals:
\begin{itemize}
\item After entering \(S\) at minute \(t\), you cannot exit \(S\) until minute \(t+6\).
\item After entering \(F\) at minute \(t\), you cannot exit \(F\) until minute \(t+4\).
\end{itemize}


\noindent\textbf{Three clockwork side-strips.}
These strips are purely mechanical and cannot be faked. A stamp at minute \(t\) means the action was performed during minute \(t\), and only someone \emph{inside} that room during minute \(t\) could perform it:
\begin{itemize}
\item The Furnace safety bell can be silenced \emph{only from inside \(F\)}; it stamped \textbf{SILENCED 21:22} and \textbf{SILENCED 21:27}.
\item The Scriptorium press can be stopped \emph{only from inside \(S\)}; it stamped \textbf{STOPPED 21:24}.
\item The Gatehouse shutter latch can be reset \emph{only from inside \(G\)}; it stamped \textbf{RESET 21:30}.
\end{itemize}

\medskip

\noindent\textbf{The locked Reliquary.}
The Reliquary is a true locked room (one door, no other exit). Its bolt stamps the minute when the door \emph{first opens} and the minute when it \emph{finally shuts}.
Once opened, a spring holds the door open for \emph{exactly three minutes}; it cannot be shut early.

A dent made the last digit hard to read; the archivist is certain it was \emph{either} \(3\) \emph{or} \(5\).
So the bolt times were \emph{either}
\[
\text{OPEN 21:23 and SHUT 21:26}
\quad\text{or}\quad
\text{OPEN 21:25 and SHUT 21:28}.
\]
A thin ash layer inside shows only \emph{one} person crossed the threshold during that single opening.

Inside, the coffer requires \emph{two full minutes} of continuous cranking to open; it cannot be completed outside the three-minute opening.

\medskip

\noindent\textbf{Ledger entries (complete for 21:15--21:35).}
\begin{itemize}
\item 21:15 --- Aegis entered \(S\).
\item 21:16 --- Athena entered \(F\).
\item 21:19 --- Quill entered \(S\).
\item 21:20 --- Kestrel entered \(G\).
\item 21:33 --- Aegis entered \(F\).
\item 21:33 --- Athena entered \(G\).
\item 21:34 --- Quill entered \(F\).
\item 21:35 --- Kestrel entered \(S\).
\end{itemize}

\medskip

\VaultDirective{Who is the culprit (thief)?}

\SolutionPage
\footnotesize

\textbf{Case 1: OPEN 21:23, SHUT 21:26.}


This window lasts three minutes, so to complete a continuous two-minute crank, the entrant must be inside \(R\) no later than 21:24.

\begin{itemize}
\item \textbf{Aegis} entered \(S\) at 21:15 and cannot exit until 21:21. The fastest route to \(R\) is \(S\to R\) (Whisper Stair) in 4 minutes, so Aegis’s earliest arrival time is 21:25. Too late.
\item \textbf{Quill} entered \(S\) at 21:19 and is sealed in until 21:25, so Quill cannot reach \(R\) by 21:24.
\item \textbf{Kestrel} entered \(G\) at 21:20; the fastest \(G\to R\) is 7 minutes, so Kestrel’s earliest arrival time is 21:27. Too late.
\item \textbf{Athena} is the only one who could, in principle, reach \(R\) by 21:24 (Athena entered \(F\) at 21:16 and can leave \(F\) from 21:20 onward; \(F\to R\) takes 3 minutes).
\end{itemize}

However, the Furnace safety bell stamped \textbf{SILENCED 21:22}, and only someone \emph{inside \(F\)} during minute 21:22 can do that. Under the ledger list, no one besides Athena can possibly be in \(F\) by 21:22.

So Athena must be in \(F\) at 21:22, which means Athena cannot leave \(F\) early enough to be in \(R\) by 21:24. Therefore no one can both satisfy the 21:22 Furnace stamp and still complete the two-minute crank within the 21:23--21:26 Reliquary opening. Thus \textbf{Case 1 is impossible}.

\medskip

\textbf{Case 2: OPEN 21:25, SHUT 21:28.}


Now to complete two continuous minutes of cranking, the entrant must be inside \(R\) by 21:26.

\begin{itemize}
\item \textbf{Quill}: sealed in \(S\) until 21:25, then needs 4 minutes to descend \(S\to R\), so earliest arrival is 21:29. Impossible.
\item \textbf{Kestrel}: from \(G\) at 21:20, earliest \(G\to R\) arrival is 21:27, leaving at most one minute before SHUT 21:28. Not enough for a two-minute crank.
\item \textbf{Athena}: Athena must be inside \(F\) at \textbf{21:27} for the second \textbf{SILENCED 21:27} stamp, so Athena cannot also be the lone person in \(R\) for two continuous minutes within 21:25--21:28.
\item \textbf{Aegis}: Aegis entered \(S\) at 21:15 and may first exit at 21:21. Taking the Whisper Stair down (4 minutes) puts Aegis at \(R\) at 21:25. Aegis can crank continuously for two full minutes (e.g., 21:25--21:27) within the three-minute opening.

After the door shuts at 21:28, taking the Ash Duct up (\(R\to F\) is 5 minutes) brings Aegis to \(F\) at 21:33, matching the ledger entry \textbf{21:33 --- Aegis entered \(F\)}.
\end{itemize}

So \textbf{Aegis is the only possible entrant} in the only possible Reliquary window.

\medskip

\textbf{Consistency snap-fit (why the rest of the stamps align).}

With Case 2 fixed and Aegis committed to \(R\) during 21:25--21:28:

\begin{itemize}
\item The Scriptorium press stamped \textbf{STOPPED 21:24}. Aegis cannot be in \(S\) at 21:24 (Aegis must have left by 21:21 to reach \(R\) at 21:25), so the only person who can be inside \(S\) at 21:24 is \textbf{Quill}, who entered \(S\) at 21:19 and is sealed in until 21:25.
\item The Gatehouse latch stamped \textbf{RESET 21:30}. Kestrel entered \(G\) at 21:20 and has no seal preventing Kestrel from remaining inside until 21:30; then \(G\to S\) takes 5 minutes, yielding Kestrel’s entry \textbf{21:35 --- Kestrel entered \(S\)}.
\item Athena must be in \(F\) at both \textbf{21:22} and \textbf{21:27}. Since Athena later entered \(G\) at 21:33 and \(F\to G\) takes 5 minutes, Athena must leave \(F\) at 21:28 (immediately after the 21:27 silencing) to arrive and enter \(G\) at 21:33.
\end{itemize}

All mechanisms agree on a single culprit.

\FinalAnswer{Aegis}

\NextPage